\section{Data, Measurement, and Empirical Results}
\subsection{Data and Measurement}
\label{sec:41:data_measurement}

% -------------------------------------------------
% 4.1.1 THE DUAL VALUATION DISTINCTION
% -------------------------------------------------
\subsubsection{The Dual Valuation Distinction}
\label{sec:411:dual_valuation}

This section constructs the capital stock series for the
choice-of-technique models of \cref{sec:3:analytical_framework}.
The BEA Fixed Assets Accounts report a single current-cost net stock.
That series treats the physical capacity of installed capital and
the monetary value of capital advanced as one magnitude, and the
conflation introduces a spurious correlation between capital
intensity and the wage share: as mechanization raises capital
intensity, Consumption of Fixed Capital (CFC) rises in the GVA
denominator, depressing GVA-basis wage shares even when the rate of
surplus-value extraction is unchanged. To eliminate this bias, I
reconstruct the stock account-by-account using the Generalized
Perpetual Inventory Method (GPIM) of \citet[App.~6.5]{Shaikh2016},
which builds capacity and profitability measures from direct nominal
investment flows and asset-specific price indexes.

\paragraph{The dual-valuation distinction.}
Estimates of capacity utilization under my unbalanced growth framework, require distinct 
capital stock variables. The capital value initially invested in plant and equipment returns
gradually to its money form as fixed assets depreciate; the
accumulated depreciation allowances remain part of total capital
advanced regardless of their monetary form, because their recovery
is what permits the reproduction of the capitalist enterprise 
\citep[App.~6.3]{Shaikh2016}. I therefore distinct two different types of
measurements of capital stocks:

\begin{enumerate}
	\item \textbf{Capital stock for capacity estimation: gross real
		GPIM stock.}
	Productive capacity depends on capital goods still physically in
	operation. The capital stock measure adequate for estimating
	capacity is the gross real stock $K^{G,R}_{j,t}$, weighted
	exclusively by a physical survival schedule. Productive capacity
	$Y^p_t$ is itself a latent variable, recovered through the
	estimation of the transformation elasticity $\theta$
	(\cref{sec:3:analytical_framework}); the gross real GPIM stock is
	the accumulation variable $k_t$ included in my estimation
	equation, not a measure of capacity itself. A machine whose book value has
	fallen to zero, but remains in service until it is physically retired, still should be accounted as physical capital in operation.  The OECD manual confirms that the gross stock, valued at ``as
	new'' prices, ``is widely used as a broad indicator of the
	productive capacity of a country'' \citep{OECD2001}.
	
	\item \textbf{Profitability measure: net current-cost stock.}
	Profitability is a monetary ratio between a profit flow and the
	value of capital advanced. The relevant measure is the net
	current-cost stock $K^{N,V}_{j,t}$, which tracks the monetary
	value of capital after financial depreciation and write-downs.
	This measure serves as the anchor for the estimation of an implicit 
	price of capital goods under a stock-flow-consistent framework but is not included within the structural estimation of non-observed productive capacities. 
\end{enumerate}

A single aggregate net stock serving both purposes is common in
empirical macroeconomics. This practice introduces a mechanical
correlation between the capital stock and the depreciation flow,
introducing potential biases to identify distributive effects on productive capacities formation. 

\paragraph{Valuation choice within the GPIM.}
The dual-valuation distinction does not require imposing a single
valuation convention throughout. \citet[App.~6.3]{Shaikh2016}
advocates the gross stock as the denominator for both the profit
rate and the capacity ratio. The GPIM he develops, however,
constructs the survival schedule $S_j(v)$, the age-price schedule
$(1-d_j)^v$, and the SFC implicit price index $P_{j,t}$ as distinct
components within a single stock-flow-consistent recursion. This
structure permits the researcher to select the valuation for each
analytical purpose without abandoning accounting coherence.

For the profit rate, the denominator must measure the monetary value
of capital advanced. The net current-cost stock serves this role: it
tracks the book value of capital that the firm must recover to
continue production. A machine that has depreciated to a low book
value may still operate at full physical capacity, but the firm's
cost-recovery problem is defined by the value it must recoup, not by
the physical state of the equipment \citep{Basu2013}. For estimates
related to capacity utilization, the capital stock input must reflect
the physical capital still in operation. The gross real GPIM stock,
weighted exclusively by the survival schedule $S_j(v)$ and excluding
the age-price factor $(1-d_j)^v$, serves this role.

The quantitative significance of the gross/net distinction is not
confined to the center. \citet{Ffrench2016} show that in Chile
depreciation absorbs between 8 and 13 percent of GDP and that net
investment is barely half of gross investment in the most recent
period. Service lives of machinery have declined from 20 years in
the 1940s to 11 years by 2005, reflecting accelerating obsolescence.
They conclude that the capital stock entering a production function
must be the gross stock of physically surviving assets, not the
depreciated monetary residual.

% -------------------------------------------------
% 4.1.2 U.S. DATA PROVENANCE AND CAPITAL RECONSTRUCTION
% -------------------------------------------------
\subsubsection{Gross Capital Stock Reconstruction and Measurement in the United States}
\label{sec:412:us_capital}

\paragraph{The Generalized Perpetual Inventory Method.}
In this part I present the procedures of construction for 
the capital stock series for the United States. Machinery (ME) and nonresidential structures (NRC) do not
share a natural physical unit; they become comparable only through a
monetary-mediated accounting procedure. I reconstruct the stock
account-by-account rather than aggregating nominal values and
deflating with a generic index. The GPIM accumulation rules for
historical-cost, current-cost, and constant-dollar stocks
\citep[App.~6.5]{Shaikh2016} are derived in detail in
Appendix~\ref{app:A1:gpim_protocol}; here I state the construction.

Physical survival of capital vintages follows a Weibull retirement
distribution. For asset class $j \in \{\text{ME}, \text{NRC}\}$,
the proportion of a vintage surviving after $v$ years is:
\begin{equation}
	S_j(v) = \exp\!\left[ -\left( \frac{v}{\lambda_j}
	\right)^{\alpha_j} \right]
	\label{eq:41:weibull_survival}
\end{equation}
where $\alpha_j$ is the Weibull shape parameter and the scale
parameter is $\lambda_j = L_j / \Gamma(1 + 1/\alpha_j)$, derived
from the mean service life $L_j$. Unlike geometric decay models that
continuously erode the physical stock, the Weibull curve ensures
that capital goods remain fully operative until they are physically
removed from service. Empirical studies of directly observed asset
discards confirm that the Weibull family provides a flexible
approximation of retirement patterns. \citet{Nomura2005}, using
Japanese micro-data on 66 asset classes, finds that almost half of
assets exhibit a progressively increasing hazard rate
($\alpha > 1$). The calibrated parameters $L_{\text{ME}} = 14$,
$\alpha_{\text{ME}} = 1.7$ and $L_{\text{NRC}} = 30$,
$\alpha_{\text{NRC}} = 1.6$ follow the BEA Fixed Assets methodology
\citep{Fraumeni1997}.

Physical survival alone cannot identify the SFC implicit price
index. The BEA current-cost net stock anchors reflect both physical
retirement and age-price decay. To recover the asset-specific price
index, I impose an externally documented geometric declining-balance
age-price profile:\footnote{The BEA's current methodology employs
	geometric depreciation, which assumes that the efficiency pattern and
	the depreciation pattern have the same form \citep{Fraumeni1997}.
	The GPIM construction separates these two dimensions: physical
	survival follows a Weibull retirement distribution, while financial
	age-price decay follows a geometric profile. This separation is
	necessary to construct distinct capacity and profitability measures.}
\begin{equation}
	V_j(v) = S_j(v) \cdot (1 - d_j)^v
	\label{eq:41:net_value_schedule}
\end{equation}
where $d_j$ is the asset-specific age-price decay rate. The
declining-balance formula applied to the calibrated service lives
yields $d_{\text{ME}} = 0.110$ and $d_{\text{NRC}} = 0.024$,
consistent with the geometric rates documented by the BEA
\citep{Fraumeni1997}.

The SFC implicit price index $P_{j,t}$ is the price level that
closes the stock-flow identity for the net current-cost stock. In
each period, the BEA anchor must satisfy
\begin{equation}
	K^{\text{NET,CC}}_{j,t}
	= I^{\text{NOM}}_{j,t}
	+ \frac{P_{j,t}}{100}
	\sum_{v=1}^{L_j} V_j(v)\, I^{\text{REAL}}_{j,t-v}
	\label{eq:41:sfc_identity}
\end{equation}
where $I^{\text{REAL}}_{j,t-v} = I^{\text{NOM}}_{j,t-v}/(P_{j,t-v}/100)$
is real investment deflated by the previously recovered SFC price.
The first term on the right-hand side is current-year nominal
investment; the second is the net-value-weighted sum of surviving
real vintages revalued at the current price. The price index is
``implicit'' because it is not directly observed: it is the unique
value that, applied to the survival-weighted real vintages,
reproduces the BEA's reported current-cost net stock. Solving
\eqref{eq:41:sfc_identity} for $P_{j,t}$:
\begin{equation}
	P_{j,t} = 100 \cdot
	\frac{K^{\text{NET,CC}}_{j,t} - I^{\text{NOM}}_{j,t}}
	{\displaystyle\sum_{v=1}^{L_j} V_j(v)\, I^{\text{REAL}}_{j,t-v}}
	\label{eq:41:sfc_price_recursion}
\end{equation}
The forward recursion proceeds period by period: given
$P_{j,t-1}, P_{j,t-2}, \ldots$, the current price $P_{j,t}$ is
uniquely determined by the requirement that the constructed stock
reproduces the BEA anchor.
Appendix~\ref{app:A1:gpim_protocol} documents the recursion, the
initial conditions, and the base-year normalization. The BEA's Fisher
quantity index is retained solely as a validation comparison and does
not enter the GPIM baseline.

The capital stock used to estimate productive capacity is the gross
real GPIM stock, constructed as the survival-weighted sum of real
investment vintages:
\begin{equation}
	K^{G,R}_{j,t} = \sum_{v=0}^{L_j} S_j(v)\,
	I^{\text{REAL}}_{j,t-v}
	\label{eq:41:gross_gpim_stock}
\end{equation}
where $I^{\text{REAL}}_{j,t-v}$ is real investment deflated by the
previously recovered SFC price. Because
\cref{eq:41:gross_gpim_stock} relies exclusively on the survival
schedule $S_j(v)$ and excludes the age-price factor $(1-d_j)^v$, it
measures the constant-price purchasing power embodied in
heterogeneous capital goods still in physical operation. Financial
write-downs, revaluation gains, and depreciation allowances do not
enter.

\paragraph{Validation and sample scope.}
The GPIM construction closes the stock-flow identity exactly: the
maximum absolute reconstruction residual across both asset classes
is $10^{-6}$ million dollars, confirming that the recovered SFC
implicit price indexes reproduce the BEA current-cost net stock
anchors. All GPIM series are normalized to $2017 = 100$.
The capital stock series has a 30-year warm-up period (1901--1930).
The GPIM stock at any year $t$ requires investment vintages
extending back $L_j$ years. Because the nominal investment series
begins in 1901, the vintage summation for the longest-lived asset
class (nonresidential structures, $L_{\text{NRC}} = 30$) is
truncated for any year before 1931. The stock values for
1901--1930 are therefore incomplete, omitting pre-1901 vintages
that would still be physically in service. I restrict the
estimation sample to 1931--2024, where the full 30-year vintage
window is populated, yielding $T = 94$ annual observations.
The sectoral scope is restricted to nonfinancial corporate (NFC)
business, excluding non-profit institutions serving households
(NPISH), government, and owner-occupied housing. This sectoral
restriction isolates the private productive business sector from
public infrastructure and imputed residential rents.

\subsubsection{Distributive Conflicts and Wage Share Dynamics in the United States}


\subsubsection{Capital Series Harmonization and Historical Sources for Chile}
\label{sec:413:chile_sfc}

This subsection describes the stock-flow consistent (SFC) methodology
used to reconstruct gross and net capital stocks for Chile,
1900--1994, in constant 2003 Chilean pesos (CLP). The procedure
enforces the accumulation identities for gross and net stocks
jointly, yields implicit depreciation and retirement rates for each
asset class, and delivers coherent series for machinery and equipment
(ME), nonresidential construction (NRC), residential construction
(RC), and their aggregates. The construction proceeds from a single
investment control total, disaggregated by asset, accumulated via
Weibull survival functions, and linked to net stocks through a
gap-ratio approach anchored to the \citet{Hofman2000b} overlap
window.

\paragraph{The stock-flow identities.}
The reconstruction is governed by two laws of motion. The gross
stock evolves according to
\begin{equation}
	\Delta K^{g}_{j,t} = I_{j,t} - R_{j,t}
	\label{eq:42:gross_identity}
\end{equation}
where $I_{j,t}$ is gross fixed investment and $R_{j,t}$ denotes
physical retirements (scrapping of obsolete vintages, demolition,
and technological withdrawal). The net stock evolves according to
\begin{equation}
	\Delta K^{n}_{j,t} = I_{j,t} - D_{j,t}
	\label{eq:42:net_identity}
\end{equation}
where $D_{j,t}$ is depreciation (the monetary consumption of fixed
capital). Combining \eqref{eq:42:gross_identity} and
\eqref{eq:42:net_identity} yields the joint SFC identity
\begin{equation}
	\Delta K^{g}_{j,t} + R_{j,t}
	= \Delta K^{n}_{j,t} + D_{j,t}
	\label{eq:42:joint_identity}
\end{equation}
which constrains the combination of depreciation and retirements at
every $t$ and for every asset $j$. Given observed stocks and
investment, the implied flows are recovered as $D_{j,t} = I_{j,t}
- \Delta K^{n}_{j,t}$ and $R_{j,t} = I_{j,t} - \Delta K^{g}_{j,t}$.
The capital gap $B_{j,t} \equiv K^{g}_{j,t} - K^{n}_{j,t}$ measures
the cumulative difference between installed capacity and book value.

\paragraph{The investment control total and the lambda bridge.}
The aggregate investment control total is drawn from the
\citet{PerezEyzaguirre2019} reconstruction of gross fixed capital
formation (GFCF), expressed in constant 2003 CLP. The control
concept is machinery plus construction, explicitly excluding
inventory variation.\footnote{Inventories are excluded from
	fixed-capital accumulation because the reconstruction models fixed
	capital only. Inventory investment does not enter the perpetual
	inventory recursion.}

Pérez-Eyzaguirre report investment in a price space and asset
aggregation that differ from the \citet{Hofman2000a} benchmark. The
two sources are linked by a level-normalization factor $\lambda$,
which corrects for definitional differences between the two series
rather than converting between price bases. The factor is estimated
as the median of the annual ratio of Hofman to Pérez investment over
the 95-year estimation window:
\begin{equation}
	\lambda = \mathrm{median}\!\left(
	\frac{I^{\text{Hofman}}_t}{I^{\text{Pérez}}_t}
	\right)_{t=1900}^{1994}
	= 4.1054
	\label{eq:42:lambda}
\end{equation}
The median rule is chosen for robustness to outliers in the early
decades, where source quality is weaker. The scaled aggregate is
then disaggregated into ME, NRC, and RC using normalized
asset-specific investment shares from \citet{Hofman2000a}:
\begin{equation}
	I_{j,t} = s_{j,t} \cdot \lambda \cdot I^{\text{Pérez}}_t,
	\quad
	\sum_{j \in \{\text{ME},\, \text{NRC},\, \text{RC}\}}
	s_{j,t} = 1
	\label{eq:42:disaggregation}
\end{equation}
for all $t$.

\paragraph{Weibull survival and gross stock accumulation.}
Physical survival of capital vintages follows a Weibull retirement
distribution. For asset class $j \in \{\text{ME}, \text{NRC},
\text{RC}\}$, the proportion of a vintage surviving after $v$ years
is
\begin{equation}
	S_j(v) = \exp\!\left[
	-\left( \frac{v}{\lambda_j} \right)^{\alpha_j}
	\right]
	\label{eq:42:weibull_chile}
\end{equation}
where $\alpha_j$ is the Weibull shape parameter and the scale
parameter is $\lambda_j = L_j / \Gamma(1 + 1/\alpha_j)$, derived
from the mean service life $L_j$. The survival parameters are
estimated from the \citet{Hofman2000b} overlap window
(1950--1994), in which both gross and net stocks are observed, by
fitting Weibull curves to the implied retirement patterns recovered
from the stock-flow identities.
Table~\ref{tab:42:survival_params} reports the calibrated
parameters.

\begin{table}[h]
	\centering
	\caption{Weibull survival parameters, Chile (1900--1994)}
	\label{tab:42:survival_params}
	\begin{tabular}{lccc}
		\toprule
		Asset & Scale ($\lambda_j$) & Shape ($\alpha_j$) & Mean life
		($L_j$, years) \\
		\midrule
		ME  & 15.90 & 2.42 & 14.1 \\
		NRC & 41.30 & 5.00 & 37.9 \\
		RC  & 51.16 & 5.00 & 47.0 \\
		\bottomrule
	\end{tabular}
\end{table}

All shape parameters exceed unity, indicating progressively
increasing hazard rates: capital goods remain fully operative in
early years and are withdrawn from service at an accelerating rate
as they approach the mean service life. The shape parameters for
NRC and RC are constrained at 5.0, reflecting the long and
relatively uniform service lives of structures. The gross stock is
accumulated as the survival-weighted sum of real investment
vintages:
\begin{equation}
	K^{g}_{j,t} = \sum_{v=0}^{L_j} S_j(v)\, I_{j,t-v}
	\label{eq:42:gross_stock_chile}
\end{equation}
Pre-1950 stocks are anchored to the observed Hofman levels at 1950
via a multiplicative factor that matches $K^{g}_{j,1950}$ by asset.

\paragraph{The gap-ratio approach and net stock construction.}
Unlike the U.S. GPIM, which recovers the net stock via a separate
geometric age-price decay rate and an SFC implicit price recursion,
the Chile reconstruction derives the net stock from the observed
gross-net relationship in the \citet{Hofman2000b} overlap window.
The capital gap ratio
\begin{equation}
	\gamma_j \equiv \frac{B_j}{K^{g}_j}
	= \frac{K^{g}_j - K^{n}_j}{K^{g}_j}
	\label{eq:42:gap_ratio}
\end{equation}
is computed by asset from the \citet{Hofman2000b} data over the
overlap period (1950--1994). Two regimes are distinguished: a
``pre'' ratio (median of years $\leq 1954$) and a ``post'' ratio
(median of years $\geq 1990$), allowing for structural shifts in
the gross-net relationship over the century. Outside the overlap
window, the gap ratio is extended at the regime-specific median.
The net stock is then recovered as
\begin{equation}
	K^{n}_{j,t} = K^{g}_{j,t} \cdot (1 - \gamma_j)
	\label{eq:42:net_stock_chile}
\end{equation}
This approach inherits the gross-net relationship observed in the
Hofman accounting reference rather than imposing an external
depreciation convention. It is appropriate for Chile because no
independent age-price decay data are available for the historical
period; the \citet{Hofman2000b} overlap provides the only direct
evidence on the joint behavior of gross and net stocks.

\paragraph{Implicit depreciation and retirement rates.}
Given the reconstructed stocks and investment flows, the implicit
annual rates are defined as
\begin{equation}
	z_{j,t} = \frac{R_{j,t}}{K^{g}_{j,t-1}},
	\qquad
	\delta_{j,t} = \frac{D_{j,t}}{K^{n}_{j,t-1}}
	\label{eq:42:implicit_rates}
\end{equation}
where $z_{j,t}$ is the retirement (depletion) rate and
$\delta_{j,t}$ is the depreciation rate. These rates are not
imposed as parameters; they are derived from the data via the
stock-flow identities
\eqref{eq:42:gross_identity}--\eqref{eq:42:net_identity}. They
serve as diagnostic indicators of capital durability and
replacement dynamics, and are used in
\cref{sec:42:us_aggregate} to validate the internal consistency of
the reconstruction.

\paragraph{Validation and estimation sample.}
The SFC construction is validated through four checks. First,
additivity: all composite stocks and flows satisfy
$K_{\text{NR}} = K_{\text{ME}} + K_{\text{NRC}}$ and
$K_{\text{T}} = K_{\text{ME}} + K_{\text{NRC}} + K_{\text{RC}}$
at every $t$. Second, non-negativity: $K^{g}_{j,t} \geq 0$,
$K^{n}_{j,t} \geq 0$, and $B_{j,t} \geq 0$ for all assets and
years. Third, rate bounds: $z_{j,t} \in [0,1]$ and
$\delta_{j,t} \in [0,1]$ for all observations from 1950 onward.
Fourth, identity residuals: $\Delta K^{g}_{j,t} - (I_{j,t} -
R_{j,t}) = 0$ and $\Delta K^{n}_{j,t} - (I_{j,t} - D_{j,t}) = 0$
to numerical precision. The estimation sample spans 1900--1994,
yielding $T = 95$ annual observations. The overlap window
(1950--1994) provides 45 years of direct \citet{Hofman2000b}
observations against which the reconstructed series are audited.

\paragraph{The BCCh forward extension (1995--2024).}
Extending the reconstruction beyond 1994 requires integrating the
official national accounts published by the Central Bank of Chile
(BCCh). A direct chain-linking of the BCCh series to the historical
baseline would introduce a level discontinuity, because the two
sources differ in base year, valuation conventions, and asset
classification. To eliminate this discontinuity while preserving the
validated 1900--1994 history, I employ a multiplicative splicing
procedure that fixes the legacy panel in place.

For the overlap year 1994, an asset-specific bridge factor is
computed as the ratio of the historical net stock to the BCCh
reported net stock:
\begin{equation}
	\beta_j = \frac{K^{n,\text{hist}}_{j,1994}}
	{K^{n,\text{BCCh}}_{j,1994}},
	\quad j \in \{\text{ME},\, \text{NRC},\, \text{RC}\}
	\label{eq:42:bridge_factor}
\end{equation}
The same factor $\beta_j$ is applied to both the BCCh net stock and
the BCCh depreciation series for all years 1995--2024. This
preserves the depreciation rates and asset dynamics implied by the
BCCh while maintaining continuity with the constant 2003 CLP
historical baseline. Table~\ref{tab:42:bridge_factors} reports the
estimated bridge factors.

\begin{table}[h]
	\centering
	\caption{Asset-specific bridge factors, 1994}
	\label{tab:42:bridge_factors}
	\begin{tabular}{lcc}
		\toprule
		Asset & Bridge factor ($\beta_j$) & Interpretation \\
		\midrule
		ME  & 5.888 & BCCh understates historical ME stock \\
		NRC & 1.033 & Near parity \\
		RC  & 1.050 & Near parity \\
		\bottomrule
	\end{tabular}
\end{table}

The large bridge factor for machinery and equipment reflects the
substantially lower level at which the BCCh reports ME stocks
relative to the Pérez-based historical reconstruction, a divergence
driven by differences in asset coverage and valuation methodology.
For nonresidential and residential construction, the bridge factors
are close to unity, indicating that the two sources agree
substantially on the level of structures.

Post-1994 real investment is inferred via the stock-flow identity:
\begin{equation}
	I_{j,t} = \Delta K^{n}_{j,t} + D_{j,t},
	\quad t \geq 1995
	\label{eq:42:post94_investment}
\end{equation}
where $D_{j,t}$ is the bridge-adjusted BCCh depreciation. The gross
stock and physical retirements for the extended period are then
reconstructed by applying the survival parameters from
Table~\ref{tab:42:survival_params} to the complete 1900--2024
investment series via
\cref{eq:42:gross_stock_chile}. This ensures that the gross stock
remains internally consistent across the splice boundary.


\subsubsection{Distributive Conflicts and Wage Share Dynamics in Chile}

\label{sec:41x:chile_distrib}
The distributive regressor in the transformation elasticity is the
wage share, $\omega_t$. For Chile, the wage share is drawn from the
revised dataset of real wages in Latin America constructed by
\citet{Astorga2017} and updated by \citet{Astorga2023}. The dataset
provides yearly series of real wages by skill level---unskilled
($w_4$), semi-skilled ($w_3$), and skilled ($w_2$)---together with
an economically-active-population (EAP) weighted average ($w$),
expressed in monthly purchasing-power-parity dollars at 1970
prices.\footnote{The revised dataset is publicly available as
	supplementary material to \citet{Astorga2023}.}

The wage share is computed as the wage-to-income ratio
$\omega_t = w_t / y_t$, where $y_t$ is real income per worker
derived from national income estimates. The functional distribution
panel extends the \citet{Astorga2023} coverage to 2024 by splicing
the revised wage series with post-2011 observations from official
Chilean statistical sources. The resulting $\omega_t$ series spans
1920--2024, yielding $T = 105$ annual observations.

\subsubsection{Aggregate Demand, Price Indices, and External Constraint variables for Chile}
\label{sec:415:chile_demand}

The aggregate demand variables---real GDP ($Y_t$), exports FOB
($X^{\text{FOB}}_t$), and gross machinery investment
($I_{\text{ME},t}$)---are drawn from the aggregate demand
reconstruction of \citet{PerezEyzaguirre2019}, expressed in constant
2003 Chilean pesos. This series revises and extends the historical
national accounts backbone of La República en Cifras
\citep{LudersDiazWagner2016}, correcting additivity and valuation
problems in the original Clio-Lab series while preserving the
sectoral GDP decomposition. Coverage spans 1860--2024.

The price and trade indices used to construct the relative capital
price proxies ($p^{\kappa}_t$), the capacity to
import ($M^{\text{cap}}_t$), are drawn from \textit{La República en Cifras}
\citep{LudersDiazWagner2016}: the real exchange rate index, the US wholesale capital price index, and the
terms of trade index. All three series are rebased to
2003 $= 100$ for internal consistency with the capital stock
series.\footnote{The terms of trade index extends back to 1810;
	the real exchange rate and US wholesale price index cover
	1900--2024.}

The capital import demand imbalance, combines sources from \textit{La República en Cifras} with my data set of harmonized capital stocks through a log difference of $log({I_{ME}}_t)$ - $log(M^{\text{cap}}_t)$, which reconciles the supply export-side identified by the index of imports capacities $(M^{\text{cap}}_t)$ of the external constraint, with the internal demand for foreign capital goods providing a both a more accurate variable for the regime shifting empirical strategy detailed below, while also a variant that allows me to check for roboustness of my estimates. 

A complementary data source for external sector variables are drawn from the World Historical Balance of Payments Database (WBOP) developed by
\citet{NievasPiketty2025}. The WBOP data set provides provides annual series on global trade flows and the world balance of payments---including goods,
services, income, and transfers---covering 57 countries over
the 1800--2025 period. Three series are used for Chile: the ratio
of manufactured goods imports to GDP (a proxy for capital goods
import dependence), net foreign income outflows as a share of GDP (a
proxy for profit remittance drain), and the net current account
balance as a share of GDP (an external balance indicator). All three
series are expressed in current US dollars at market exchange rates
and are available at the WBOP online repository.\footnote{All WBOP
	series are available at \texttt{wbop.world} and are also hosted in
	the World Inequality Database (\texttt{wid.world}), together with a
	detailed replication package including raw data sources, methods,
	and codes.}




%%%%%%%%%%%%%%%%%%%%%%%%%%%%%%%%%%%%%%%%%%%%%%%%%%%%%%%
%%%%% 4.2. Econometric Strategy IM-OLS Cointegration
%%%%%%%%%%%%%%%%%%%%%%%%%%%%%%%%%%%%%%%%%%%%%%%%%%%%%%%

\subsection{Econometric Strategy: Integrated Modified OLS and Cointegration}



%%%%%%%%%%%%%%%%%%%%%%%%%%%%%%%%%%%%%%%%%%%%%%%%%
%%%%% 4.3. US Results Aggregate Capital
%%%%%%%%%%%%%%%%%%%%%%%%%%%%%%%%%%%%%%%%%%%%%%%%%


\subsection{United States Empirical Evidence: Aggregate Capital Accumulation and Capacity Utilization}

\subsubsection{Baseline Empirical Estimates}

\subsubsection{Transformation Elasticity and Capacity Utilization Estimates}


%%%%%%%%%%%%%%%%%%%%%%%%%%%%%%%%%%%%%%%%%%%%%%%%%
%%%%% 4.4. US Results Aggregate Capital
%%%%%%%%%%%%%%%%%%%%%%%%%%%%%%%%%%%%%%%%%%%%%%%%%


\subsection{United States Empirical Evidence: Capital Heterogeneity and the Mechanization Gap}


\subsubsection{Heterogeneous Capital Stocks Empirical Estimates}


\subsubsection{Aggregate Transformation Elasticity and Capacity Utilization Estimates with Heterogeneous Capital Stocks}


%%%%%%%%%%%%%%%%%%%%%%%%%%%%%%%%%%%%%%%%%%%%%%%%%
%%%%% 4.5. Chile Results 
%%%%%%%%%%%%%%%%%%%%%%%%%%%%%%%%%%%%%%%%%%%%%%%%%
\subsection{Chilean Empirical Evidence: The Peripheral and External Constraints}

\subsubsection{Econometric Strategy: Threshold Conditional Mixed Process Estimation}


\subsubsection{Empirical Estimates of the Balance-of-Payments Constraint}

\subsubsection{Transformation and Capacity Utilization und Financial Subordination}



%%%%%%%%%%%%%%%%%%%%%%%%%%%%%%%%%%%%%%%%%%%%%%%%%
%%%%% 4.6 Comparative and Historical Analysis 
%%%%%%%%%%%%%%%%%%%%%%%%%%%%%%%%%%%%%%%%%%%%%%%%%
\subsection{Comparative Analysis and Theoretical Synthesis}


